\let\negmedspace\undefined
\let\negthickspace\undefined
\documentclass[journal]{article}
\usepackage[a5paper, margin=10mm, onecolumn]{geometry}
\usepackage{lmodern} % Ensure lmodern is loaded for pdflatex

\setlength{\headheight}{1cm} % Set the height of the header box
\setlength{\headsep}{0mm}     % Set the distance between the header box and the top of the text

\usepackage{gvv-book}
\usepackage{gvv}
\usepackage{cite}
\usepackage{textcomp}
\usepackage{amsmath,amssymb,amsfonts,amsthm}
\usepackage{algorithmic}
\usepackage{graphicx}
\graphicspath{{./figs/}}
\usepackage{textcomp}
\usepackage{xcolor}
\usepackage{txfonts}
\usepackage{listings}
\usepackage{enumitem}
\usepackage{mathtools}
\usepackage{gensymb}
\usepackage{comment}
\usepackage[breaklinks=true]{hyperref}
\usepackage{tkz-euclide} 
\usepackage{listings}
\usepackage{gvv}                                        
\def\inputGnumericTable{}                                 
\usepackage[latin1]{inputenc}                                
\usepackage{color}                                            
\usepackage{array}                                            
\usepackage{longtable}                                       
\usepackage{calc}                                             
\usepackage{multirow}                                         
\usepackage{hhline}                                           
\usepackage{ifthen}                                           
\usepackage{lscape}
\usepackage{circuitikz}
\tikzstyle{block} = [rectangle, draw, fill=blue!20, 
text width=4em, text centered, rounded corners, minimum height=3em]
\tikzstyle{sum} = [draw, fill=blue!10, circle, minimum size=1cm, node distance=1.5cm]
\tikzstyle{input} = [coordinate]
\tikzstyle{output} = [coordinate]


\begin{document}
	
	\bibliographystyle{IEEEtran}
	\vspace{3cm}
	
\title{9.5.10}
\author{EE25BTECH11047 - RAVULA SHASHANK REDDY}
\maketitle
\hrulefill
\bigskip 

\renewcommand{\thetable}{\theenumi}
\setlength{\intextsep}{10pt}

\textbf{Question:} \\

If the sum of zeroes of the polynomial is $\sqrt{2}$, then find the value of $k$ .
\begin{align*}
    p(x)=2x^2-k\sqrt{2}x+1
\end{align*}

\textbf{Solution:}
\begin{align}
\vec{x}^\top \vec{V} \vec{x} + 2\vec{u}^\top \vec{x} + f = 0\\
\vec{x} = \vec{h} + \kappa \vec{m}.
\end{align}

Substituting the line into the conic gives
\begin{align}
(\vec{h} + \kappa \vec{m})^\top \vec{V} (\vec{h} + \kappa \vec{m}) + 2\vec{u}^\top (\vec{h} + \kappa \vec{m}) + f &= 0 \\
\Rightarrow\;
\kappa^2(\vec{m}^\top \vec{V} \vec{m}) + 2\kappa\, \vec{m}^\top(\vec{V}\vec{h} + \vec{u}) + (\vec{h}^\top \vec{V} \vec{h} + 2\vec{u}^\top \vec{h} + f) &= 0\\
\kappa_1 + \kappa_2 = -\frac{2\,\vec{m}^\top(\vec{V}\vec{h}+\vec{u})}{\vec{m}^\top \vec{V} \vec{m}}.
\end{align}

Given:(Roots are given by intersection of conic with x-axis)
\begin{align}
\vec{V} &=
\myvec{
2 & 0\\[4pt]
0 & 0
}, \quad
\vec{u} =
\myvec{
-\dfrac{k\sqrt{2}}{2}\\[4pt]
0
}, \quad
f = 1\\
\vec{h}&=\myvec{0 \\ 0},\quad \vec{m}=\myvec{1\\0}\\
\vec{V}\vec{h} + \vec{u} &=
\myvec{
- \dfrac{k\sqrt{2}}{2}\\[4pt]
0
},\\[4pt]
\vec{m}^\top \vec{V} \vec{m} &= 2.\\
\vec{m}^\top (\vec{V}\vec{h} + \vec{u})&= - \dfrac{k\sqrt{2}}{2}
\end{align}

\begin{align}
\kappa_1 + \kappa_2 
&= -\frac{2(- \dfrac{k\sqrt{2}}{2})}{2}=\sqrt{2}.
\end{align}
\begin{align}
\Rightarrow\;\boxed{k = 2}.
\end{align}
\newpage
\begin{figure}
    \centering
    \includegraphics[width=1.0\linewidth]{figs/Fig_1.png}
    \caption{Caption}
    \label{fig:placeholder}
\end{figure}
\end{document}

